\documentclass{prettydoc}

% --- Document metadata -----------------------------------------------------
% \accentword{} inside the title colours a word (blue by default; pass a colour
% for the optional argument, e.g. \accentword[brandpink]{...}).
\title{Working \accentword{together}, simply.}
\kicker{A quick orientation}
\lead{A short, brand-styled template you can drop any public-facing content
into — briefs, guides, proposals or client notes. Replace the copy; the
styling looks after itself.}
% Optional extra metadata — unused fields simply don't render.

\doctype{Brief}
\revision{A}
\documentId{FD-DC-LTX-10001}  % shown in the footer; omit to fall back to the build ID
\author{Fiddlie}
\date{June 2026}              % the meta-block year (2026) is pulled from this
% \website{} defaults to fiddlie.com (last line of the top-right meta block)

\begin{document}
\maketitle

% --- An eyebrow label introduces a section, just like the cover -----------
\eyebrow{What this template gives you}

The building blocks below are the same ones used across Fiddlie's
public-facing documents. Each is a single command, so a document stays
readable in source and consistent in print.

% --- Numbered cards: a pale-blue chip above a bold heading ----------------
% \card auto-numbers (01, 02, ...). Pair each with whatever content you like.
\card{A branded cover}

\begin{itemize}
  \item A full-bleed \textbf{tricolour bar} along the top edge
  \item The \texttt{fiddlie.} wordmark with a right-aligned metadata block
  \item A two-tone display title and an optional lead paragraph
\end{itemize}

\card{Eyebrow labels}

\begin{itemize}
  \item Use \verb|\eyebrow{...}| to introduce a section in letterspaced blue caps
  \item The same treatment appears above the cover title via \verb|\kicker{}|
\end{itemize}

\card{Numbered cards}

\begin{itemize}
  \item \verb|\card{Heading}| draws the chip and heading you see here
  \item Numbering is automatic; pass \verb|\card[A]{...}| to set a label, or
        \verb|\card*{...}| to reuse the current number
\end{itemize}

\card{Blue-dot lists}

\begin{itemize}
  \item Ordinary \texttt{itemize} lists pick up the brand bullet automatically
  \item \textbf{Inline emphasis} works exactly as usual
\end{itemize}

% --- The next part opens on a fresh page with a full display heading -------
\clearpage
\eyebrow{Putting it to use}

\section{Write \accentword[brandpink]{ordinary} LaTeX.}

There is no special body layout to learn. Headings, paragraphs and lists all
inherit the brand styling, so you can focus on the content.

\subsection{Section headings}

Use \verb|\section{...}| for big display headings — they accept \verb|\accentword{}|
for a two-tone effect. Use \verb|\subsection{...}| for the compact bold lead-ins
that sit above a short run of bullets, like this one.

\subsection{Two-tone accents}

\begin{itemize}
  \item \verb|\accentword{word}| colours a word in brand \accentword{blue}
  \item \verb|\accentword[brandpink]{word}| uses \accentword[brandpink]{pink}
  \item \verb|\accentword[brandyellow]{word}| uses \accentword[brandyellow]{yellow}
\end{itemize}

\subsection{Footer and metadata}

Every page carries the document reference — \verb|\documentId{}| with the
revision — bottom-left, with the build ID on a line beneath it (a LOCAL
marker until a release build sets it). A small bold \accentword{01 / 02}
page indicator sits squared into the bottom-right corner. The cover's meta
block is driven by \verb|\doctype{}|, \verb|\revision{}| and \verb|\website{}|;
the year is pulled from \verb|\date{}|.

% --- A closing sign-off ----------------------------------------------------
\closing{Swap in your own content and the document stays on-brand —
predictable structure, so the writing has room to breathe.}

\end{document}
